\section{Source and Delivery Note}
For an adult parish assembly on the Twentieth Sunday after Pentecost,
11 October 2026, in the 1962 Roman Missal. This homily follows the
expansive study's reading \emph{The word heals and teaches faith},
with the compatible relation between lament and hope developed in
\emph{Learning thanksgiving in exile}. Its decisive detail is the
father's departure in John~4:50 before the servants' confirmation in
4:51--53. The matching hour reveals mercy already given during the
unseen journey, and the household's belief makes its fruit communal.
The Epistle's wise conduct and mutual service, and the prayers for
pardon, medicine and obedience, join this growth in trust and shared faith.

The speech contains 1,377 spoken words. At 115--125 words per minute,
its estimated duration is 11.0--12.0 minutes. A silent editorial read-through
checked sentence sense, transitions and oral clarity; no audible rehearsal
or timed human delivery is claimed. This note and the references are
outside the speech. The homily applies in either commemoration branch:
it makes no claim that the conditional Maternity prayers are said at
every Mass. It is authored preaching prose, not an official liturgical text.

\section{References}
\begin{description}[style=nextline,leftmargin=0pt,labelsep=0pt,itemsep=.35em]
\item[Appointed Scripture and prayers]
John~4:46b--53; Ephesians~5:15--21; Psalms~136:1 (137) and
118:49--50 (119). Biblical quotations follow Douay--Rheims--Challoner,
Project Gutenberg eBook 1581. \work{Missale Romanum}, Vatican, 1962,
Twentieth Sunday after Pentecost, pp.~412--413: Collect, Secret and
Postcommunion. The prayers are explained, not supplied as substitute
translations or recited prayers. The full texts and their historical
witnesses belong to the companion expansive study.

\item[Gregory the Great]
\work{Homilies on the Gospels}, II.28.1, on the ruler's faith and
Christ's power without bodily travel. Latin Wikisource transcription,
contributors including Mizardellorsa, CC BY-SA 3.0:
\url{https://la.wikisource.org/wiki/Homiliarum_in_Evangelia}.
The homily paraphrases Gregory's argument. It does not attribute to
Augustine Gregory's account of the ruler's initial faith.

\item[John Chrysostom and Augustine]
Chrysostom, \work{Homilies on Ephesians}, XIX, exposition of 5:15--17
and 5:21, NPNF I.13: prudent fidelity and reciprocal service.
Augustine, \work{Expositions of the Psalms}, Ps.~118, \S\S50--51
(NPNF I.8, heading Ps.~119): remembrance as fulfilment of a promise
and comfort for Christ's body under trial, reaching to resurrection.

\item[Liturgical reception]
Ildefonso Schuster, \work{The Sacramentary}, Arthur Levelis-Marke's
translation, III (1927), ``Twentieth Sunday after Pentecost,''
pp.~175--178: pardon and peaceful service, Eucharistic nourishment
and medicine, hope and obedient reception. Historical biblical and
NPNF English and Schuster's text retain their recorded United States
public-domain status; the project license covers the original exposition.
\end{description}
